% !PRINTED 57
\weberpage{57}

%<<CONT
der analytischen Geometrie) der \gesperrt{Flächeninhalt des Periodenparallelogramms}.

Durch Aneinanderreihen kongruenter Periodenparallelogramme
läßt sich die ganze $u$-Ebene einfach und lückenlos überdecken.
Entsprechende Punkte zweier verschiedener dieser Parallelogramme
sind die Repräsentanten von $u$-Werten, die sich um ganzzahlige
Vielfache der Perioden unterscheiden, und die \gesperrt{nach dem Modul
$(\omega_1,\omega_2)$ kongruent} oder, wenn kein Zweifel möglich ist, kurzweg
\gesperrt{kongruent} genannt werden. Das Zeichen der Kongruenz ist

\[ u \equiv u' \ (\mathrm{mod}\ \omega_1,\ \omega_2)\,, \]

und besagt, daß $u'$ die Form $u+m_1\omega_1+m_2\omega_2$ hat, wenn $m_1$, $m_2$
ganze Zahlen sind.

% Source figure: Weber_Algebra_Bd3.pdf, PDF page 79, printed p.57, Fig.1; exact crop including caption and every original label.
% Former placeholder description retained: Vier aneinandergrenzende Periodenparallelogramme im Gitternetz der Perioden $\omega_1,\omega_2$, mit Eckpunkten $u_0$, $u_0+\omega_1$, $u_0+2\omega_1$, $u_0+\omega_2$, $u_0+\omega_1+\omega_2$, $u_0+2\omega_1+\omega_2$, $u_0+2\omega_2$, $u_0+\omega_1+2\omega_2$, $u_0+2\omega_1+2\omega_2$, und den inneren (gestrichelt verbundenen) Punkten $u_1$, $u_2$, $u_3$, $u_4$.
\weberfigureimage{b3_i0078_fig1_source_20260831.png}{0.66}

So stellt die Fig.\,1 vier aneinander gelagerte Periodenparallelogramme
 dar; die Punkte $u_1$, $u_2$, $u_3$, $u_4$ sind kongruent:

\[ u_2 = u_1+\omega_1\,,\quad u_3 = u_1+\omega_1+\omega_2\,,\quad u_4 = u_1+\omega_2\,, \]

und $u_1$, $u_2$, $u_3$, $u_4$ bilden ebenfalls die Ecken eines Periodenparallelogramms.

\gesperrt{Die doppelt periodische Funktion $\varphi(u)$ hat in allen
kongruenten Punkten denselben Wert}, und der ganze Wertvorrat
 dieser Funktion erschöpft sich also \gesperrt{in einem Periodenparallelogramm}, wenn von zwei gegenüberliegenden Seiten nur
\gesperrt{die eine} mit zum Parallelogramm gerechnet wird.

Wir ziehen aus unseren allgemeinen Voraussetzungen die
Folgerung:

1. \gesperrt{Es existiert (außer der Konstanten) keine doppelt
periodische Funktion, die im Periodenparallelogramm
frei von Unstetigkeitspunkten ist;} denn eine solche Funk-
%CONT>>
